\documentclass[11pt,letterpaper]{article}

% ============================================================
% COSMIC UNIVERSALISM WEBSITE IMPLEMENTATION MANUAL
% Compile with: pdfLaTeX
% Verified baseline date: July 18, 2026
% ============================================================

% ------------------------------------------------------------
% Core packages
% ------------------------------------------------------------
\usepackage[margin=0.78in,headheight=16pt]{geometry}
\usepackage[T1]{fontenc}
\usepackage[utf8]{inputenc}
\usepackage{lmodern}
\usepackage[scaled=0.94]{helvet}
\usepackage{microtype}
\usepackage{parskip}
\usepackage{setspace}
\usepackage{enumitem}
\usepackage{booktabs}
\usepackage{tabularx}
\usepackage{array}
\usepackage{longtable}
\usepackage{xcolor}
\usepackage{graphicx}
\usepackage{fancyhdr}
\usepackage{lastpage}
\usepackage{titlesec}
\usepackage{hyperref}
\usepackage{bookmark}
\usepackage{listings}
\usepackage[most]{tcolorbox}
\usepackage{amssymb}
\usepackage{pifont}

% ------------------------------------------------------------
% Document typography
% ------------------------------------------------------------
\renewcommand{\familydefault}{\sfdefault}
\setstretch{1.08}
\setlength{\parindent}{0pt}
\setlength{\parskip}{0.65em}

% ------------------------------------------------------------
% Cosmic Universalism color palette
% ------------------------------------------------------------
\definecolor{CUNavy}{HTML}{050812}
\definecolor{CUNavyTwo}{HTML}{080D1A}
\definecolor{CUSurface}{HTML}{0C1424}
\definecolor{CUGold}{HTML}{DDBF72}
\definecolor{CUGoldLight}{HTML}{FFF8DD}
\definecolor{CUCyan}{HTML}{35CDEB}
\definecolor{CUCyanLight}{HTML}{78E8FF}
\definecolor{CUBlue}{HTML}{668ED1}
\definecolor{CUBlueLight}{HTML}{9EB9EA}
\definecolor{CUAmber}{HTML}{D99A32}
\definecolor{CURed}{HTML}{B84C4C}
\definecolor{CUGreen}{HTML}{2C8A65}
\definecolor{CULight}{HTML}{F3F6FA}
\definecolor{CUBorder}{HTML}{CCD5E0}
\definecolor{CUText}{HTML}{18212E}
\definecolor{CUMuted}{HTML}{5C6878}
\definecolor{CUCodeBackground}{HTML}{F6F8FB}
\definecolor{CUCodeBorder}{HTML}{BFCADA}

% ------------------------------------------------------------
% Project metadata
% ------------------------------------------------------------
\newcommand{\ProjectTitle}{Cosmic Universalism Website}
\newcommand{\DocumentTitle}{Local Website Implementation Manual}
\newcommand{\RepositoryName}{CosmicUniversalismStatement}
\newcommand{\RepositoryURL}{https://github.com/willmaddock/CosmicUniversalismStatement}
\newcommand{\WebsiteBranch}{website}
\newcommand{\WebsiteDirectory}{website/}
\newcommand{\DocumentVersion}{1.0}
\newcommand{\VerificationDate}{July 18, 2026}
\newcommand{\AuthorName}{William Maddock}

% ------------------------------------------------------------
% Hyperlinks and PDF metadata
% ------------------------------------------------------------
\hypersetup{
    colorlinks=true,
    linkcolor=CUBlue,
    urlcolor=CUBlue,
    citecolor=CUBlue,
    pdftitle={Cosmic Universalism Website Implementation Manual},
    pdfauthor={William Maddock},
    pdfsubject={Astro website implementation instructions},
    pdfkeywords={Cosmic Universalism, Astro, TypeScript, GitHub Desktop,
                 PyCharm, ChatGPT Desktop, Codex, GitHub Pages}
}

% ------------------------------------------------------------
% Header and footer
% ------------------------------------------------------------
\pagestyle{fancy}
\fancyhf{}
\fancyhead[L]{\small\textbf{\ProjectTitle}}
\fancyhead[R]{\small\DocumentTitle}
\fancyfoot[L]{\small Version \DocumentVersion}
\fancyfoot[C]{\small \thepage\ of \pageref{LastPage}}
\fancyfoot[R]{\small Branch: \texttt{\WebsiteBranch}}
\renewcommand{\headrulewidth}{0.5pt}
\renewcommand{\footrulewidth}{0.5pt}

% ------------------------------------------------------------
% Section formatting
% ------------------------------------------------------------
\titleformat{\section}
  {\Large\bfseries\color{CUNavy}}
  {\thesection}
  {0.8em}
  {}
  [\vspace{0.2em}\color{CUCyan}\titlerule]

\titleformat{\subsection}
  {\large\bfseries\color{CUNavyTwo}}
  {\thesubsection}
  {0.7em}
  {}

\titleformat{\subsubsection}
  {\normalsize\bfseries\color{CUBlue}}
  {\thesubsubsection}
  {0.6em}
  {}

\titlespacing*{\section}{0pt}{2.2em}{1em}
\titlespacing*{\subsection}{0pt}{1.6em}{0.7em}
\titlespacing*{\subsubsection}{0pt}{1.2em}{0.4em}

% ------------------------------------------------------------
% Lists
% ------------------------------------------------------------
\setlist[itemize]{
    leftmargin=1.6em,
    itemsep=0.3em,
    topsep=0.3em
}

\setlist[enumerate]{
    leftmargin=2em,
    itemsep=0.45em,
    topsep=0.4em
}

\newcommand{\checkbox}{\(\square\)}
\newcommand{\checked}{\ding{51}}

% ------------------------------------------------------------
% Table helpers
% ------------------------------------------------------------
\newcolumntype{Y}{>{\raggedright\arraybackslash}X}
\newcolumntype{L}[1]{>{\raggedright\arraybackslash}p{#1}}

% ------------------------------------------------------------
% General information boxes
% ------------------------------------------------------------
\newtcolorbox{purposebox}{
    enhanced,
    breakable,
    colback=CULight,
    colframe=CUBlue,
    boxrule=0.8pt,
    arc=2mm,
    left=3mm,
    right=3mm,
    top=2.5mm,
    bottom=2.5mm,
    title=\textbf{Purpose},
    coltitle=white,
    colbacktitle=CUBlue
}

\newtcolorbox{checkpointbox}{
    enhanced,
    breakable,
    colback=CUCyan!6,
    colframe=CUCyan,
    boxrule=0.8pt,
    arc=2mm,
    left=3mm,
    right=3mm,
    top=2.5mm,
    bottom=2.5mm,
    title=\textbf{Checkpoint},
    coltitle=CUNavy,
    colbacktitle=CUCyanLight
}

\newtcolorbox{safetybox}{
    enhanced,
    breakable,
    colback=CUAmber!8,
    colframe=CUAmber,
    boxrule=0.8pt,
    arc=2mm,
    left=3mm,
    right=3mm,
    top=2.5mm,
    bottom=2.5mm,
    title=\textbf{Safety Rule},
    coltitle=white,
    colbacktitle=CUAmber
}

\newtcolorbox{warningbox}{
    enhanced,
    breakable,
    colback=CURed!6,
    colframe=CURed,
    boxrule=0.8pt,
    arc=2mm,
    left=3mm,
    right=3mm,
    top=2.5mm,
    bottom=2.5mm,
    title=\textbf{Do Not Continue Until Resolved},
    coltitle=white,
    colbacktitle=CURed
}

\newtcolorbox{successbox}{
    enhanced,
    breakable,
    colback=CUGreen!7,
    colframe=CUGreen,
    boxrule=0.8pt,
    arc=2mm,
    left=3mm,
    right=3mm,
    top=2.5mm,
    bottom=2.5mm,
    title=\textbf{Successful Result},
    coltitle=white,
    colbacktitle=CUGreen
}

\newtcolorbox{notebox}{
    enhanced,
    breakable,
    colback=CUCodeBackground,
    colframe=CUBorder,
    boxrule=0.7pt,
    arc=2mm,
    left=3mm,
    right=3mm,
    top=2.5mm,
    bottom=2.5mm
}

% ------------------------------------------------------------
% Listing styles
% ------------------------------------------------------------
\lstdefinestyle{terminalstyle}{
    basicstyle=\ttfamily\small,
    backgroundcolor=\color{CUNavyTwo},
    keywordstyle=\color{CUCyanLight},
    stringstyle=\color{CUGoldLight},
    commentstyle=\color{CUBlueLight},
    identifierstyle=\color{white},
    showstringspaces=false,
    columns=fullflexible,
    keepspaces=true,
    breaklines=true,
    breakatwhitespace=false,
    frame=none,
    tabsize=2,
    upquote=true
}

\lstdefinestyle{promptstyle}{
    basicstyle=\ttfamily\small,
    backgroundcolor=\color{CUCodeBackground},
    keywordstyle=\color{CUBlue},
    stringstyle=\color{CUGreen},
    commentstyle=\color{CUMuted},
    identifierstyle=\color{CUText},
    showstringspaces=false,
    columns=fullflexible,
    keepspaces=true,
    breaklines=true,
    breakatwhitespace=false,
    frame=none,
    tabsize=2,
    upquote=true
}

\lstdefinestyle{codestyle}{
    basicstyle=\ttfamily\small,
    backgroundcolor=\color{CUCodeBackground},
    keywordstyle=\color{CUBlue},
    stringstyle=\color{CUGreen},
    commentstyle=\color{CUMuted},
    identifierstyle=\color{CUText},
    showstringspaces=false,
    columns=fullflexible,
    keepspaces=true,
    breaklines=true,
    breakatwhitespace=false,
    frame=none,
    tabsize=2,
    upquote=true
}

% ------------------------------------------------------------
% Copy-and-paste boxes
% ------------------------------------------------------------
\newtcblisting{commandbox}[1]{
    enhanced,
    breakable,
    listing only,
    listing engine=listings,
    listing options={style=terminalstyle},
    colback=CUNavyTwo,
    colframe=CUCyan,
    coltitle=CUNavy,
    colbacktitle=CUCyanLight,
    boxrule=0.8pt,
    arc=2mm,
    left=2mm,
    right=2mm,
    top=1mm,
    bottom=1mm,
    title=\textbf{#1}
}

\newtcblisting{promptbox}[1]{
    enhanced,
    breakable,
    listing only,
    listing engine=listings,
    listing options={style=promptstyle},
    colback=CUCodeBackground,
    colframe=CUBlue,
    coltitle=white,
    colbacktitle=CUBlue,
    boxrule=0.8pt,
    arc=2mm,
    left=2mm,
    right=2mm,
    top=1mm,
    bottom=1mm,
    title=\textbf{Copy-and-Paste Prompt: #1}
}

\newtcblisting{codebox}[1]{
    enhanced,
    breakable,
    listing only,
    listing engine=listings,
    listing options={style=codestyle},
    colback=CUCodeBackground,
    colframe=CUCodeBorder,
    coltitle=white,
    colbacktitle=CUNavyTwo,
    boxrule=0.8pt,
    arc=2mm,
    left=2mm,
    right=2mm,
    top=1mm,
    bottom=1mm,
    title=\textbf{#1}
}

% ------------------------------------------------------------
% Convenience commands
% ------------------------------------------------------------
\newcommand{\filepath}[1]{\texttt{#1}}
\newcommand{\menu}[1]{\textbf{\textsf{#1}}}
\newcommand{\branchname}[1]{\texttt{#1}}
\newcommand{\statusline}[2]{\textbf{#1:} #2}

% ============================================================
% DOCUMENT
% ============================================================
\begin{document}

% ------------------------------------------------------------
% Cover page
% ------------------------------------------------------------
\begin{titlepage}
\pagecolor{CUNavy}
\color{white}
\thispagestyle{empty}

\vspace*{0.8in}

{\small\bfseries\color{CUCyanLight}
COSMIC UNIVERSALISM COMPUTATIONAL INTELLIGENCE INITIATIVE}

\vspace{0.45in}

{\Huge\bfseries
\ProjectTitle}

\vspace{0.2in}

{\LARGE\color{CUGoldLight}
\DocumentTitle}

\vspace{0.4in}

{\large
A controlled, step-by-step workflow for building a local Astro website
with ChatGPT Desktop, PyCharm, GitHub Desktop, and GitHub Pages.}

\vfill

\begin{tcolorbox}[
    colback=CUNavyTwo,
    colframe=CUBlue,
    boxrule=0.8pt,
    arc=2mm,
    width=\textwidth
]
\color{white}
\begin{tabularx}{\textwidth}{>{\bfseries}L{1.65in}Y}
Repository & \RepositoryName \\
Development branch & \texttt{\WebsiteBranch} \\
Website directory & \texttt{\WebsiteDirectory} \\
Primary framework & Astro with TypeScript \\
Editor & PyCharm with the Astro plugin \\
Version control & GitHub Desktop \\
Planning environment & ChatGPT Work \\
Implementation environment & ChatGPT Codex \\
Document version & \DocumentVersion \\
Technical baseline verified & \VerificationDate \\
\end{tabularx}
\end{tcolorbox}

\vfill

{\large\bfseries \AuthorName}

\vspace{0.1in}

{\small
\href{\RepositoryURL}{\RepositoryURL}}

\vspace{0.5in}

\end{titlepage}

\pagecolor{white}
\color{CUText}

% ------------------------------------------------------------
% Document control
% ------------------------------------------------------------
\section*{Document Control}
\addcontentsline{toc}{section}{Document Control}

\begin{tabularx}{\textwidth}{>{\bfseries}L{1.8in}Y}
\toprule
Document & \ProjectTitle: \DocumentTitle \\
Version & \DocumentVersion \\
Status & Initial controlled implementation guide \\
Repository & \href{\RepositoryURL}{\RepositoryName} \\
Protected branch & \branchname{main} \\
Development branch & \branchname{\WebsiteBranch} \\
Website source directory & \filepath{\WebsiteDirectory} \\
Intended platform & macOS with GitHub Desktop and PyCharm \\
Deployment target & GitHub Pages \\
Last technical verification & \VerificationDate \\
\bottomrule
\end{tabularx}

\begin{notebox}
Software interfaces and action versions can change. Before the final
deployment, compare the deployment section with the current official
Astro and GitHub documentation. The local build and branch-protection
principles remain valid even when interface labels change.
\end{notebox}

\clearpage
\tableofcontents
\clearpage

% ============================================================
\section{Purpose and Intended Outcome}
% ============================================================

\begin{purposebox}
This manual guides the creation of a local, production-quality website
for the \textbf{Cosmic Universalism Statement} repository while
protecting the existing research, converter, media, and main branch.
\end{purposebox}

The intended repository arrangement is:

\begin{codebox}{Target repository structure}
CosmicUniversalismStatement/
|-- .github/
|-- ResearchFiles/
|-- cosmic_converter/
|-- Images/
|-- media/
|-- website/
|   |-- public/
|   |-- src/
|   |-- astro.config.mjs
|   |-- package.json
|   |-- package-lock.json
|   `-- tsconfig.json
|-- README.md
`-- LICENSE.md
\end{codebox}

The branch arrangement is:

\begin{codebox}{Branch strategy}
main
`-- Canonical Cosmic Universalism research repository

website
`-- Website source, website content, and website deployment workflow
\end{codebox}

\begin{safetybox}
The \branchname{website} branch is a separate version of the complete
repository. Work performed on that branch does not alter
\branchname{main} unless changes are deliberately merged into
\branchname{main}.
\end{safetybox}

% ============================================================
\section{Application Responsibilities}
% ============================================================

\begin{tabularx}{\textwidth}{L{1.55in}Y}
\toprule
\textbf{Application} & \textbf{Primary responsibility} \\
\midrule
GitHub Desktop &
Create and switch branches, review file changes, create commits,
fetch, push, and publish branches. \\

PyCharm &
Inspect and edit Astro, TypeScript, HTML, CSS, Markdown, SVG, and
configuration files. Run the local terminal and inspect project files. \\

ChatGPT Work &
Plan phases, structure research, write specifications, organize
content, review design decisions, and prepare implementation tasks. \\

ChatGPT Codex &
Inspect the repository, create or edit code, run commands, execute
builds, report errors, and provide file-by-file change summaries. \\

Browser &
View the website at the local Astro development address and test
desktop and mobile behavior. \\

Astro &
Generate the static production website. \\

GitHub Pages &
Host the generated public website after the local build is approved. \\
\bottomrule
\end{tabularx}

\begin{checkpointbox}
Each application has one primary responsibility. Avoid asking several
applications to manage Git simultaneously. Use GitHub Desktop as the
main visual source-control interface.
\end{checkpointbox}

% ============================================================
\section{Phase 1: Preflight Checklist}
% ============================================================

Complete this checklist before creating website files.

\begin{itemize}
    \item[\checkbox] GitHub Desktop is installed and signed in.
    \item[\checkbox] The repository is available locally.
    \item[\checkbox] PyCharm is installed.
    \item[\checkbox] The JetBrains Astro plugin is installed.
    \item[\checkbox] Node.js 22.12.0 or newer is installed.
    \item[\checkbox] Node.js 24 LTS is preferred when available.
    \item[\checkbox] npm is available.
    \item[\checkbox] ChatGPT Desktop is installed and signed in.
    \item[\checkbox] The approved homepage mockup is saved.
    \item[\checkbox] The website blueprint is saved.
\end{itemize}

\subsection{Verify Node.js and npm}

Open Terminal or PyCharm's built-in terminal and run:

\begin{commandbox}{Verify Node.js and npm}
node --version
npm --version
\end{commandbox}

A suitable Node.js result is:

\begin{codebox}{Expected Node.js result}
v22.12.0 or newer
\end{codebox}

A Node 24 result may resemble:

\begin{codebox}{Recommended example}
v24.x.x
\end{codebox}

\begin{warningbox}
Do not continue if the terminal reports
\texttt{command not found: node} or if the Node version is older than
22.12.0.
\end{warningbox}

% ============================================================
\section{Phase 2: Protect the Repository with a Website Branch}
% ============================================================

\subsection{Fetch the current repository}

In GitHub Desktop:

\begin{enumerate}
    \item Select the repository \textbf{\RepositoryName}.
    \item Confirm the current branch is \branchname{main}.
    \item Click \menu{Fetch origin}.
    \item Allow GitHub Desktop to finish synchronizing.
\end{enumerate}

\subsection{Create the website branch}

In GitHub Desktop:

\begin{enumerate}
    \item Click \menu{Current Branch}.
    \item Click \menu{New Branch}.
    \item Enter the branch name:
\end{enumerate}

\begin{codebox}{Branch name}
website
\end{codebox}

\begin{enumerate}[resume]
    \item Set \branchname{main} as the base branch.
    \item Click \menu{Create Branch}.
    \item Click \menu{Publish branch}.
\end{enumerate}

\begin{checkpointbox}
The top of GitHub Desktop must display:

\begin{center}
\texttt{Current Branch: website}
\end{center}

Do not create or edit website files while GitHub Desktop displays
\branchname{main}.
\end{checkpointbox}

\subsection{Verify the branch in Terminal}

Open Terminal in the repository root and run:

\begin{commandbox}{Verify the active branch}
git branch --show-current
\end{commandbox}

The exact result should be:

\begin{successbox}
\begin{center}
\texttt{website}
\end{center}
\end{successbox}

% ============================================================
\section{Phase 3: Configure PyCharm}
% ============================================================

\subsection{Open the complete repository}

In PyCharm:

\begin{enumerate}
    \item Select \menu{File \textrightarrow{} Open}.
    \item Select the complete \filepath{CosmicUniversalismStatement}
          folder.
    \item Do not select only \filepath{ResearchFiles/}.
    \item Do not select only \filepath{cosmic\_converter/}.
    \item Allow PyCharm to finish indexing the project.
\end{enumerate}

\subsection{Install the Astro plugin}

In PyCharm:

\begin{enumerate}
    \item Open \menu{PyCharm \textrightarrow{} Settings}.
    \item Select \menu{Plugins}.
    \item Select \menu{Marketplace}.
    \item Search for \texttt{Astro}.
    \item Install the official Astro plugin.
    \item Restart PyCharm if requested.
\end{enumerate}

Then inspect:

\begin{codebox}{Astro language-server location}
Settings
  -> Languages & Frameworks
  -> TypeScript
  -> Astro
\end{codebox}

Enable Astro support if that option is available.

\subsection{Configure Node.js in PyCharm}

Open:

\begin{codebox}{PyCharm Node configuration}
Settings
  -> Languages & Frameworks
  -> Node.js
\end{codebox}

Confirm:

\begin{itemize}
    \item Node interpreter: Node 22.12.0 or newer.
    \item Package manager: npm.
\end{itemize}

\subsection{Open the built-in terminal}

In PyCharm:

\begin{codebox}{Open Terminal}
View
  -> Tool Windows
  -> Terminal
\end{codebox}

Run:

\begin{commandbox}{Final PyCharm environment check}
git branch --show-current
node --version
npm --version
\end{commandbox}

Record the results:

\begin{longtable}{L{2.2in}L{3.8in}}
\toprule
\textbf{Check} & \textbf{Observed result} \\
\midrule
Active Git branch & \\
Node.js version & \\
npm version & \\
Astro plugin installed & Yes / No \\
Repository opened at root & Yes / No \\
\bottomrule
\end{longtable}

% ============================================================
\section{Phase 4: Create the ChatGPT Project}
% ============================================================

\subsection{Create the project}

In ChatGPT Desktop:

\begin{enumerate}
    \item Create a new Project.
    \item Name it:
\end{enumerate}

\begin{codebox}{ChatGPT Project name}
Cosmic Universalism Website
\end{codebox}

\subsection{Add the permanent project instructions}

Open the Project's instruction field and paste the following.

\begin{promptbox}{Project Instructions}
PROJECT: Cosmic Universalism Website

REPOSITORY:
willmaddock/CosmicUniversalismStatement

LOCAL REPOSITORY:
CosmicUniversalismStatement

DEVELOPMENT BRANCH:
website

WEBSITE DIRECTORY:
website/

TECHNOLOGY:
Astro, TypeScript, HTML, modern CSS, Markdown, SVG,
and lightweight browser JavaScript.

DEVELOPMENT APPLICATIONS:
- PyCharm for reviewing and editing source files
- GitHub Desktop for branches, commits, and pushes
- ChatGPT Work for planning, research organization,
  design, and content
- ChatGPT Codex for repository inspection, coding,
  commands, builds, and tests

PROTECTION RULES:
1. Never modify the main branch.
2. Confirm the current branch is website before making changes.
3. Keep new website source code inside website/.
4. Do not delete or rewrite existing repository files unless
   explicitly directed.
5. Do not modify the CU-Time converter mathematics without a
   separate technical review.
6. Do not commit node_modules/, dist/, .DS_Store, private keys,
   tokens, passwords, or secret .env files.
7. Make small, reviewable changes.
8. Explain planned file changes before implementing them.
9. Run npm run build after major changes.
10. Do not deploy until the local production build succeeds.
11. Do not commit, push, merge, or deploy unless explicitly asked.

CU CONTENT RULES:
1. Preserve sub-ctom and ctom as different terms.
2. Do not place the Observable Universe marker inside sub-ctom.
3. Keep scientific references separate from CU theoretical
   propositions.
4. Clearly label open problems, contradictions, and superseded
   interpretations.
5. Preserve the 3.108-trillion-year full Cosmic Breath framework.
6. Do not silently change established project terminology.

VISUAL DIRECTION:
A realistic combination of NASA scientific visualization,
modern academic publishing, cinematic cosmic documentary,
and restrained philosophical symbolism.

COLORS:
- Deep navy-black background
- White-gold for the Great Baking Will and sub-ztom
- Cyan for empirical references and the Observable Universe
- Translucent blue for compression structures
- Amber for open questions

WORKFLOW:
Plan in Work.
Implement and test in Codex.
Review files in PyCharm.
Review and commit changes in GitHub Desktop.
\end{promptbox}

\subsection{Add reference materials}

Add these reference materials to the Project:

\begin{itemize}
    \item The approved Cosmic Universalism homepage mockup.
    \item The website blueprint or outline.
    \item Any canonical research documents needed for the first pages.
\end{itemize}

Use recognizable filenames:

\begin{codebox}{Suggested reference filenames}
cosmic-universalism-homepage-reference.png
cosmic-universalism-website-blueprint.md
\end{codebox}

\begin{safetybox}
Do not upload passwords, GitHub tokens, private SSH keys, secret
environment files, or unrelated personal documents.
\end{safetybox}

% ============================================================
\section{Phase 5: Connect ChatGPT Desktop and PyCharm}
% ============================================================

\subsection{Enable Work with Apps}

In ChatGPT Desktop:

\begin{enumerate}
    \item Open \menu{ChatGPT \textrightarrow{} Settings}.
    \item Find \menu{Work with Apps}.
    \item Enable the feature.
    \item Open \menu{Manage Apps}.
    \item Confirm whether PyCharm appears as a supported application.
    \item Grant only the permissions required for the project.
\end{enumerate}

The exact interface wording can vary by desktop-app release.

\subsection{Prepare the PyCharm window}

Before attaching PyCharm:

\begin{itemize}
    \item Keep only the Cosmic Universalism repository open.
    \item Close unrelated private projects.
    \item Close files containing credentials or secrets.
    \item Confirm the active branch is \branchname{website}.
\end{itemize}

\subsection{Use reviewable changes}

For the first implementation phases:

\begin{itemize}
    \item Ask for a plan before edits.
    \item Ask for a list of files that will change.
    \item Review proposed diffs.
    \item Avoid automatic application of broad edits.
    \item Test each approved change locally.
\end{itemize}

% ============================================================
\section{Phase 6: Create the Initial Work Plan}
% ============================================================

Open ChatGPT Work inside the Cosmic Universalism Website Project.

Paste:

\begin{promptbox}{Phase 1 Planning}
We are beginning the Cosmic Universalism website project.

Use the Project Instructions, website blueprint, and approved
homepage reference image.

For this task, do not edit any files.

Create a concise Phase 1 implementation plan for a local Astro
website. The first milestone must include only:

1. Astro project initialization inside website/
2. Global design tokens
3. Base page layout
4. Header navigation
5. Static homepage hero
6. Placeholder Cosmic Breath SVG
7. Three summary cards
8. Research classification legend

For each item, identify:

- the goal
- the expected folders and files
- dependencies
- accessibility requirements
- mobile behavior
- testing requirements
- completion criteria

Do not implement anything yet.
Do not modify the repository.
Do not commit, push, merge, or deploy.
\end{promptbox}

Review the plan before proceeding.

\begin{checkpointbox}
The plan should keep all new application files inside
\filepath{website/}. It should not propose changing the converter
mathematics or rewriting the research repository.
\end{checkpointbox}

% ============================================================
\section{Phase 7: Inspect the Repository with Codex}
% ============================================================

Switch to ChatGPT Codex and open the local repository.

Paste:

\begin{promptbox}{Repository Inspection}
Inspect this repository without changing any files.

Confirm:

1. The repository root
2. The current Git branch
3. Whether a website/ directory already exists
4. Whether package.json exists at the repository root
5. Whether Node.js and npm are available
6. The detected Node.js and npm versions
7. Whether there are uncommitted changes
8. Which GitHub Pages workflow files currently exist
9. Which branches those workflows trigger on
10. Whether any existing workflow could deploy from the
    website branch

Do not install packages.
Do not create files.
Do not modify files.
Do not switch branches.
Do not commit.
Do not push.
Do not merge.
Do not deploy.

Report the findings first.
\end{promptbox}

\begin{warningbox}
If Codex reports that the current branch is \branchname{main}, stop.
Return to GitHub Desktop and switch to \branchname{website}.
\end{warningbox}

% ============================================================
\section{Phase 8: Initialize Astro}
% ============================================================

\subsection{Preferred Codex initialization prompt}

After the inspection confirms the correct branch, paste:

\begin{promptbox}{Initialize Astro}
Initialize a new Astro website inside the website/ directory.

Requirements:

- Use the minimal Astro template.
- Use strict TypeScript.
- Install dependencies.
- Use npm as the package manager.
- Do not initialize a nested Git repository.
- Do not modify files outside website/.
- Do not modify the existing GitHub Pages workflow.
- Do not deploy.
- Do not commit or push.
- After creation, run npm run build.
- Report every created file.
- Report all warnings or errors.
- Confirm whether a nested website/.git directory exists.

Stop after the initial successful Astro build.
\end{promptbox}

\subsection{Manual terminal alternative}

To perform initialization manually from the repository root:

\begin{commandbox}{Launch the Astro setup wizard}
npm create astro@latest website
\end{commandbox}

Choose options equivalent to:

\begin{codebox}{Astro wizard selections}
Template: Minimal
Install dependencies: Yes
TypeScript: Strict
Initialize a new Git repository: No
\end{codebox}

The wizard wording may vary slightly.

\subsection{Expected project structure}

After initialization, PyCharm should show:

\begin{codebox}{Expected Astro structure}
website/
|-- public/
|-- src/
|   `-- pages/
|       `-- index.astro
|-- .gitignore
|-- astro.config.mjs
|-- package.json
|-- package-lock.json
|-- README.md
`-- tsconfig.json
\end{codebox}

\begin{warningbox}
There should not be a nested \filepath{website/.git/} directory.
A nested Git repository would interfere with the parent repository's
source control.
\end{warningbox}

% ============================================================
\section{Phase 9: Run and Test the Initial Website}
% ============================================================

\subsection{Start the development server}

In PyCharm's terminal:

\begin{commandbox}{Run the local development server}
cd website
npm run dev
\end{commandbox}

Astro normally reports an address similar to:

\begin{codebox}{Local development address}
http://localhost:4321/
\end{codebox}

Open the address in Safari or Chrome.

To stop the server, press:

\begin{codebox}{Stop the server}
Control + C
\end{codebox}

\subsection{Run the production build}

From \filepath{website/}:

\begin{commandbox}{Build the production website}
npm run build
\end{commandbox}

A successful build creates:

\begin{codebox}{Generated build directory}
website/dist/
\end{codebox}

\subsection{Preview the production build}

\begin{commandbox}{Preview the production build}
npm run preview
\end{commandbox}

\begin{successbox}
The initial milestone is successful when:

\begin{itemize}
    \item The development server launches.
    \item The local page opens in a browser.
    \item \texttt{npm run build} completes without errors.
    \item \texttt{npm run preview} displays the production build.
\end{itemize}
\end{successbox}

% ============================================================
\section{Phase 10: Review and Commit the Initial Astro Project}
% ============================================================

\subsection{Inspect the Git changes}

Return to GitHub Desktop.

The changes should be concentrated inside:

\begin{codebox}{Expected change location}
website/
\end{codebox}

Do not commit:

\begin{itemize}
    \item \filepath{node\_modules/}
    \item \filepath{dist/}
    \item \filepath{.DS\_Store}
    \item secret \filepath{.env} files
    \item access tokens
    \item passwords
    \item private keys
\end{itemize}

Confirm that \filepath{website/.gitignore} includes:

\begin{codebox}{Minimum ignore entries}
node_modules
dist
\end{codebox}

\subsection{Create the first commit}

Use this commit summary:

\begin{codebox}{Commit message}
Create initial Astro website
\end{codebox}

In GitHub Desktop:

\begin{enumerate}
    \item Confirm the button says \menu{Commit to website}.
    \item Review every selected file.
    \item Click \menu{Commit to website}.
    \item Click \menu{Push origin}.
\end{enumerate}

\begin{warningbox}
Do not click the commit button if it says
\menu{Commit to main}.
\end{warningbox}

% ============================================================
\section{Phase 11: Establish the Website Architecture}
% ============================================================

Use ChatGPT Work to prepare the architecture before implementation.

\begin{promptbox}{Website Architecture Plan}
Create the implementation specification for the initial Cosmic
Universalism website architecture.

Do not edit files.

The application root is website/.

Plan the following structure:

website/
  public/
    diagrams/
    images/
    media/
    tools/
  src/
    components/
    content/
      framework/
      research/
      glossary/
      media/
      archive/
    layouts/
    pages/
      research/
    styles/

Plan these initial components:

- Header.astro
- Hero.astro
- CosmicBreathDiagram.astro
- BreathSummary.astro
- ResearchLegend.astro
- Footer.astro

Plan these initial layouts:

- BaseLayout.astro
- ResearchLayout.astro

Plan these initial pages:

- index.astro
- framework.astro
- cosmic-breath.astro
- cu-time.astro
- research/index.astro
- media.astro
- about.astro

Plan these initial stylesheets:

- tokens.css
- global.css
- typography.css

For every file, explain:

- its responsibility
- what it imports
- what content it owns
- what content it must not duplicate
- accessibility requirements
- responsive behavior
- completion criteria

Do not implement the files.
Do not modify the repository.
Do not commit, push, merge, or deploy.
\end{promptbox}

% ============================================================
\section{Phase 12: Implement the Design Foundation}
% ============================================================

\subsection{Planning prompt}

\begin{promptbox}{Design System Specification}
Using the approved Cosmic Universalism mockup and Project
Instructions, prepare a detailed design-system specification.

Do not edit files.

Define:

1. Color tokens
2. Typography tokens
3. Spacing scale
4. Maximum content width
5. Research reading width
6. Border and surface styles
7. Button styles
8. Focus states
9. Research classification styles
10. Reduced-motion behavior
11. Desktop breakpoints
12. Mobile breakpoints

Use these visual principles:

- deep navy-black background
- white-gold for the Great Baking Will and sub-ztom
- cyan for empirical references and the Observable Universe
- translucent blue for compression structures
- amber for open questions
- high readability
- restrained glow
- no video-game HUD styling
- no cryptocurrency styling
- no excessive particle effects

Return the specification only.
Do not implement it.
\end{promptbox}

\subsection{Implementation prompt}

After approving the specification, switch to Codex and paste:

\begin{promptbox}{Implement Design Tokens and Base Layout}
Implement only the approved design tokens and base page layout.

Requirements:

- Work only inside website/.
- Create maintainable CSS custom properties.
- Create a semantic BaseLayout.astro.
- Add accessible metadata defaults.
- Add a skip-to-content link.
- Add visible keyboard focus styles.
- Add reduced-motion support.
- Do not build the final header or hero yet.
- Do not add heavy dependencies.
- Do not modify existing research or converter files.
- Run npm run build.
- Report all changed files.
- Report warnings and errors.
- Do not commit, push, merge, or deploy.
\end{promptbox}

Suggested commit message after review:

\begin{codebox}{Commit message}
Add global CU design system and base layout
\end{codebox}

% ============================================================
\section{Phase 13: Build the Header and Static Hero}
% ============================================================

\subsection{Work planning prompt}

\begin{promptbox}{Header and Hero Specification}
Using the approved homepage mockup, create an implementation
specification for the homepage header and static hero only.

Do not edit files.

Specify:

- exact component names
- semantic HTML structure
- navigation labels
- accessible mobile navigation behavior
- hero heading hierarchy
- exact hero text
- button labels
- provisional button destinations
- responsive desktop and mobile layout
- placeholder behavior for the Cosmic Breath SVG
- keyboard focus behavior
- reduced-motion requirements
- acceptance criteria

Required navigation labels:

Framework
Cosmic Breath
CU-Time
Research
Media
About
GitHub

Required site identity:

COSMIC UNIVERSALISM
Computational Intelligence Initiative

Required hero heading:

COSMIC UNIVERSALISM

Required hero description:

A philosophical and computational framework exploring
consciousness, time, existence, and the full Cosmic Breath.

Required buttons:

Explore the Framework
Open the CU Statement

Required research-status line:

Scientific references distinguished from CU theoretical
propositions.

Do not implement the specification.
\end{promptbox}

\subsection{Codex implementation prompt}

\begin{promptbox}{Implement Header and Static Hero}
Implement only the approved header and static homepage hero.

Requirements:

- Work only inside website/.
- Use semantic and accessible HTML.
- Use the existing design tokens.
- Keep content separate from decorative styling.
- Implement a responsive desktop and mobile arrangement.
- Use a static placeholder for the Cosmic Breath diagram.
- Do not add animation yet.
- Do not add a JavaScript framework.
- Do not modify research files.
- Do not modify converter mathematics.
- Run npm run build.
- Provide a file-by-file change summary.
- Report warnings and errors.
- Do not commit, push, merge, or deploy.
\end{promptbox}

Suggested commit message:

\begin{codebox}{Commit message}
Add homepage header and static hero
\end{codebox}

% ============================================================
\section{Phase 14: Build the Cosmic Breath SVG}
% ============================================================

\subsection{Planning prompt}

\begin{promptbox}{Cosmic Breath SVG Specification}
Create a technical specification for a responsive, accessible,
static SVG visualization of the Cosmic Breath.

Do not edit files.

The SVG must include these exact labels:

- sub-ztom
- atom
- btom
- ctom
- ztom

The center must contain a white-gold sub-ztom seed.

The Observable Universe marker must:

- be cyan
- appear inside the ctom region
- display "You Are Here"
- display "Observable Universe"
- not appear inside sub-ctom
- not overlap the central seed

The diagram should look like a realistic scientific website
visualization built from SVG shapes, gradients, strokes, masks,
and labels.

It must not resemble:

- a gaming HUD
- a magical portal
- a cryptocurrency graphic
- a crowded science-fiction dashboard

Specify:

- SVG viewBox
- grouping strategy
- shell geometry
- text-label strategy
- responsive scaling
- accessibility title and description
- mobile fallback
- reduced-motion behavior for later animation
- acceptance criteria

Do not implement the SVG yet.
\end{promptbox}

\subsection{Implementation prompt}

\begin{promptbox}{Implement Static Cosmic Breath SVG}
Implement the approved static Cosmic Breath SVG.

Requirements:

- Work only inside website/.
- Use an Astro component named CosmicBreathDiagram.astro.
- Preserve the exact approved tom labels.
- Place the cyan Observable Universe marker inside ctom.
- Do not introduce a sub-ctom label unless the specification
  explicitly requires it.
- Add an accessible SVG title and description.
- Ensure text remains readable at supported desktop widths.
- Provide a simplified readable mobile arrangement.
- Do not animate the SVG yet.
- Run npm run build.
- Report changed files, warnings, and errors.
- Do not commit, push, merge, or deploy.
\end{promptbox}

Suggested commit message:

\begin{codebox}{Commit message}
Add static Cosmic Breath SVG visualization
\end{codebox}

% ============================================================
\section{Phase 15: Build the Summary Cards and Research Legend}
% ============================================================

\subsection{Required card content}

\begin{codebox}{Expansion card}
Expansion Phase
2.8 trillion years
sub-ztom -> atom
\end{codebox}

\begin{codebox}{Compression card}
Compression Phase
308 billion years
btom -> ztom
\end{codebox}

\begin{codebox}{Full Breath card}
Full Cosmic Breath
3.108 trillion years
Expansion + Compression
\end{codebox}

\subsection{Required research classifications}

\begin{codebox}{Research legend}
Empirical Reference
CU Mathematical Model
CU Theoretical Proposition
Open Question
\end{codebox}

\subsection{Implementation prompt}

\begin{promptbox}{Implement Summary Cards and Research Legend}
Implement the approved Cosmic Breath summary cards and research
classification legend.

Requirements:

- Work only inside website/.
- Create BreathSummary.astro.
- Create ResearchLegend.astro.
- Use semantic HTML.
- Use the approved design tokens.
- Stack cards vertically on narrow screens.
- Preserve all numerical values and labels exactly.
- Use cyan for Empirical Reference.
- Use white-gold for CU Mathematical Model.
- Use translucent blue for CU Theoretical Proposition.
- Use amber for Open Question.
- Do not add unsupported scientific claims.
- Do not add animation.
- Run npm run build.
- Provide a file-by-file change summary.
- Do not commit, push, merge, or deploy.
\end{promptbox}

Suggested commit message:

\begin{codebox}{Commit message}
Add Cosmic Breath summary and research legend
\end{codebox}

% ============================================================
\section{Phase 16: Repeatable Feature Workflow}
% ============================================================

Use the following sequence for every feature:

\begin{enumerate}
    \item Plan one small feature in ChatGPT Work.
    \item Review the plan.
    \item Implement only that feature in Codex.
    \item Review every changed file in PyCharm.
    \item Run the development server.
    \item Inspect the feature in the browser.
    \item Test desktop and mobile widths.
    \item Run the production build.
    \item Review GitHub Desktop changes.
    \item Commit to \branchname{website}.
    \item Push origin.
\end{enumerate}

\subsection{Generic planning prompt}

\begin{promptbox}{Generic Feature Planning}
Plan the following website feature:

[INSERT FEATURE NAME]

Do not edit files.

Repository rules:

- Current development branch: website
- Website application directory: website/
- Do not modify main
- Do not modify existing research or converter files unless
  explicitly directed
- Do not change CU-Time mathematics
- Preserve CU terminology exactly

Provide:

1. Goal
2. User experience
3. Files to create
4. Files to modify
5. Semantic HTML structure
6. Accessibility requirements
7. Responsive behavior
8. Data or content requirements
9. Testing steps
10. Completion criteria
11. Risks or unresolved questions

Do not implement.
Do not commit, push, merge, or deploy.
\end{promptbox}

\subsection{Generic implementation prompt}

\begin{promptbox}{Generic Feature Implementation}
Implement only this approved feature:

[INSERT FEATURE NAME]

Approved specification:

[PASTE OR REFERENCE THE APPROVED SPECIFICATION]

Rules:

- Confirm the current branch is website before editing.
- Work only inside website/ unless the approved specification
  explicitly identifies another file.
- Make the smallest maintainable change.
- Use semantic HTML.
- Preserve accessibility.
- Preserve responsive behavior.
- Do not add unnecessary dependencies.
- Do not modify CU-Time mathematics.
- Do not alter research conclusions.
- Run npm run build after editing.
- Report every changed file.
- Report all warnings and errors.
- Do not commit, push, merge, or deploy.
\end{promptbox}

\subsection{Generic review prompt}

\begin{promptbox}{Review the Completed Feature}
Review the current uncommitted website changes.

Do not edit files during the first review.

Check:

1. Whether all changes are on the website branch
2. Whether files outside website/ changed
3. Whether the implementation matches the approved specification
4. Accessibility
5. Responsive behavior
6. Broken links
7. Incorrect paths
8. Incorrect CU terminology
9. sub-ctom versus ctom confusion
10. Observable Universe marker placement
11. Build warnings
12. Unnecessary dependencies
13. Secret or generated files that should not be committed
14. Whether npm run build succeeds

Report findings from highest to lowest priority.
Do not commit, push, merge, or deploy.
\end{promptbox}

% ============================================================
\section{Phase 17: Recommended Feature Order}
% ============================================================

Develop features in this order:

\begin{enumerate}
    \item Initial Astro application
    \item Global design tokens
    \item Base layout
    \item Header
    \item Static hero
    \item Static Cosmic Breath SVG
    \item Summary cards
    \item Research-classification legend
    \item Mobile navigation
    \item Footer
    \item Framework overview page
    \item CU Statement page
    \item Cosmic Breath page
    \item Research index
    \item Research article layout
    \item Glossary
    \item CU-Time converter wrapper
    \item Media page
    \item Subtle motion
    \item Reduced-motion review
    \item Accessibility review
    \item Production deployment
\end{enumerate}

\begin{safetybox}
Do not begin complex animation, WebGL, search, a database, a Go
backend, or user accounts during the initial website milestone.
\end{safetybox}

% ============================================================
\section{Phase 18: Local Quality Checks}
% ============================================================

Run these commands from \filepath{website/}:

\begin{commandbox}{Development server}
npm run dev
\end{commandbox}

\begin{commandbox}{Production build}
npm run build
\end{commandbox}

\begin{commandbox}{Production preview}
npm run preview
\end{commandbox}

Use this checklist:

\begin{itemize}
    \item[\checkbox] Homepage loads.
    \item[\checkbox] No browser-console errors appear.
    \item[\checkbox] All navigation items are visible.
    \item[\checkbox] Keyboard focus is visible.
    \item[\checkbox] The skip link works.
    \item[\checkbox] The hero remains readable.
    \item[\checkbox] The SVG labels are correct.
    \item[\checkbox] The Observable Universe marker is inside ctom.
    \item[\checkbox] Summary values are correct.
    \item[\checkbox] Research classifications are distinguishable.
    \item[\checkbox] The page works at mobile width.
    \item[\checkbox] The page works at desktop width.
    \item[\checkbox] Reduced-motion settings are respected.
    \item[\checkbox] \texttt{npm run build} succeeds.
\end{itemize}

% ============================================================
\section{Phase 19: Commit Discipline}
% ============================================================

Use small, descriptive commits.

Recommended sequence:

\begin{codebox}{Suggested commit history}
Create initial Astro website
Add global CU design system and base layout
Add homepage header and static hero
Add static Cosmic Breath SVG visualization
Add Cosmic Breath summary and research legend
Add responsive mobile navigation
Add Framework overview page
Add CU Statement page
Add research library structure
Integrate CU-Time converter page
Add reduced-motion support
Prepare website branch deployment
\end{codebox}

Before every commit:

\begin{itemize}
    \item[\checkbox] GitHub Desktop displays \branchname{website}.
    \item[\checkbox] The changed-file list has been reviewed.
    \item[\checkbox] No secret files are selected.
    \item[\checkbox] No \filepath{node\_modules/} files are selected.
    \item[\checkbox] No \filepath{dist/} files are selected.
    \item[\checkbox] The production build succeeds.
    \item[\checkbox] The commit message describes one logical change.
\end{itemize}

% ============================================================
\section{Phase 20: GitHub Pages Deployment}
% ============================================================

\begin{warningbox}
Do not perform this phase until the local production build has been
reviewed and approved.

The repository may already contain a GitHub Pages workflow. Inspect
all existing workflows before adding another deployment workflow.
If an existing workflow deploys from \branchname{main}, future pushes
to \branchname{main} could overwrite a site deployed from
\branchname{website}.
\end{warningbox}

\subsection{Configure Astro for the repository path}

Edit \filepath{website/astro.config.mjs}:

\begin{codebox}{astro.config.mjs}
import { defineConfig } from 'astro/config';

export default defineConfig({
  site: 'https://willmaddock.github.io',
  base: '/CosmicUniversalismStatement',
  output: 'static',
});
\end{codebox}

Because the site uses a repository subpath, manually written internal
links must account for the configured base path.

\subsection{Deployment-review prompt}

Before creating the workflow, use:

\begin{promptbox}{Deployment Safety Review}
Inspect the repository's current GitHub Actions workflows.

Do not edit files.

Determine:

1. Which workflows deploy GitHub Pages
2. Which branches trigger each workflow
3. Whether a main-branch push could overwrite a deployment from
   the website branch
4. Whether GitHub Pages is currently configured for Actions
5. Whether the Astro project path should be ./website
6. Whether astro.config.mjs has the correct site and base values
7. Whether package-lock.json is committed
8. Whether all internal links support the configured base path

Recommend the safest deployment arrangement for keeping:

- main as the canonical research branch
- website as the permanent website branch

Do not modify workflows.
Do not commit, push, merge, or deploy.
\end{promptbox}

\subsection{Website-branch deployment workflow}

As of the verification date shown on this document, the official Astro
documentation uses:

\begin{itemize}
    \item \texttt{actions/checkout@v7}
    \item \texttt{withastro/action@v6}
    \item \texttt{actions/deploy-pages@v5}
\end{itemize}

Create this file only after reviewing existing workflows:

\begin{codebox}{.github/workflows/deploy-website.yml}
name: Deploy Cosmic Universalism Website

on:
  push:
    branches: [website]
  workflow_dispatch:

permissions:
  contents: read
  pages: write
  id-token: write

concurrency:
  group: pages
  cancel-in-progress: false

jobs:
  build:
    runs-on: ubuntu-latest

    steps:
      - name: Checkout repository
        uses: actions/checkout@v7

      - name: Install, build, and upload Astro site
        uses: withastro/action@v6
        with:
          path: ./website
          node-version: 24

  deploy:
    needs: build
    runs-on: ubuntu-latest

    environment:
      name: github-pages
      url: ${{ steps.deployment.outputs.page_url }}

    steps:
      - name: Deploy to GitHub Pages
        id: deployment
        uses: actions/deploy-pages@v5
\end{codebox}

\begin{notebox}
GitHub Action major versions can change. Recheck the official Astro
GitHub Pages documentation before committing this workflow.
\end{notebox}

\subsection{Configure GitHub Pages}

On GitHub:

\begin{codebox}{GitHub Pages settings}
Repository
  -> Settings
  -> Pages
  -> Source
  -> GitHub Actions
\end{codebox}

The expected public address is:

\begin{codebox}{Expected public URL}
https://willmaddock.github.io/CosmicUniversalismStatement/
\end{codebox}

\subsection{Deployment implementation prompt}

\begin{promptbox}{Prepare Deployment}
Prepare the approved GitHub Pages deployment for the Astro website.

Requirements:

- Confirm the current branch is website.
- Inspect existing workflow files before editing.
- Do not allow an old workflow to unexpectedly overwrite the
  website deployment.
- Configure website/astro.config.mjs with:
  site: https://willmaddock.github.io
  base: /CosmicUniversalismStatement
  output: static
- Create or update only the approved website deployment workflow.
- Use ./website as the Astro project path.
- Run npm run build locally.
- Check internal links and asset paths against the base path.
- Report every changed file.
- Do not commit or push.
- Do not trigger a deployment.
\end{promptbox}

% ============================================================
\section{Phase 21: Troubleshooting}
% ============================================================

\subsection{Node command not found}

Symptom:

\begin{codebox}{Error}
zsh: command not found: node
\end{codebox}

Response:

\begin{itemize}
    \item Install a supported Node.js release.
    \item Close and reopen Terminal.
    \item Run \texttt{node --version}.
\end{itemize}

\subsection{Wrong branch}

Symptom:

\begin{codebox}{Unexpected result}
main
\end{codebox}

Response:

\begin{enumerate}
    \item Stop editing.
    \item Save or discard uncommitted changes appropriately.
    \item Open GitHub Desktop.
    \item Switch to \branchname{website}.
    \item Run \texttt{git branch --show-current} again.
\end{enumerate}

\subsection{Nested Git repository}

Symptom:

\begin{codebox}{Unwanted directory}
website/.git/
\end{codebox}

Response:

\begin{itemize}
    \item Do not commit yet.
    \item Remove only the nested \filepath{website/.git/} directory.
    \item Do not remove the parent repository's \filepath{.git/}.
    \item Recheck GitHub Desktop.
\end{itemize}

\subsection{Port 4321 is already in use}

Run Astro on another port:

\begin{commandbox}{Alternative local port}
npm run dev -- --port 4322
\end{commandbox}

Then open:

\begin{codebox}{Alternative address}
http://localhost:4322/
\end{codebox}

\subsection{PyCharm does not recognize Astro files}

Check:

\begin{itemize}
    \item The Astro plugin is installed and enabled.
    \item PyCharm was restarted after installation.
    \item The complete repository is open.
    \item Node.js is configured.
    \item Astro language-server support is enabled.
\end{itemize}

\subsection{GitHub Pages returns a 404}

Check:

\begin{itemize}
    \item The workflow succeeded.
    \item Pages source is set to GitHub Actions.
    \item \texttt{site} is correct.
    \item \texttt{base} is
          \texttt{/CosmicUniversalismStatement}.
    \item Internal links include the repository base path.
    \item Static asset paths are compatible with the base path.
\end{itemize}

\subsection{The build fails}

Run:

\begin{commandbox}{Clean dependency installation}
cd website
rm -rf node_modules
npm install
npm run build
\end{commandbox}

Do not delete \filepath{package-lock.json} unless dependency recovery
specifically requires regenerating it.

% ============================================================
\section{Phase 22: Initial Project Checkpoint}
% ============================================================

Complete only these tasks before beginning homepage implementation:

\begin{itemize}
    \item[\checkbox] Fetch origin in GitHub Desktop.
    \item[\checkbox] Create the \branchname{website} branch.
    \item[\checkbox] Publish the \branchname{website} branch.
    \item[\checkbox] Verify the branch in Terminal.
    \item[\checkbox] Verify Node.js.
    \item[\checkbox] Verify npm.
    \item[\checkbox] Open the repository root in PyCharm.
    \item[\checkbox] Install the Astro plugin.
    \item[\checkbox] Create the ChatGPT Project.
    \item[\checkbox] Add the Project Instructions.
    \item[\checkbox] Add the homepage reference image.
    \item[\checkbox] Add the website blueprint.
    \item[\checkbox] Run the planning prompt.
    \item[\checkbox] Run the no-edit repository inspection.
\end{itemize}

Record the checkpoint:

\begin{longtable}{L{2.5in}L{3.5in}}
\toprule
\textbf{Checkpoint item} & \textbf{Result} \\
\midrule
Current branch & \\
Node.js version & \\
npm version & \\
PyCharm Astro plugin & \\
ChatGPT Project created & \\
Reference mockup added & \\
Blueprint added & \\
Repository inspection complete & \\
Uncommitted changes before setup & \\
Existing deployment workflows & \\
\bottomrule
\end{longtable}

% ============================================================
\section{Official Reference Links}
% ============================================================

These links were used to establish the technical baseline of this
manual.

\begin{itemize}
    \item Astro installation and Node.js requirements:

    \href{https://docs.astro.build/en/install-and-setup/}
    {https://docs.astro.build/en/install-and-setup/}

    \item Astro editor and JetBrains setup:

    \href{https://docs.astro.build/en/editor-setup/}
    {https://docs.astro.build/en/editor-setup/}

    \item Astro GitHub Pages deployment:

    \href{https://docs.astro.build/en/guides/deploy/github/}
    {https://docs.astro.build/en/guides/deploy/github/}

    \item GitHub Desktop branch management:

    \href{https://docs.github.com/en/desktop/making-changes-in-a-branch/managing-branches-in-github-desktop}
    {GitHub Desktop branch-management documentation}

    \item ChatGPT Projects:

    \href{https://help.openai.com/en/articles/10169521-projects-in-chatgpt}
    {OpenAI Projects documentation}

    \item Work with Apps on macOS:

    \href{https://help.openai.com/en/articles/10119604-work-with-apps-on-macos}
    {OpenAI Work with Apps documentation}

    \item ChatGPT Work and Codex:

    \href{https://help.openai.com/en/articles/20001275-chatgpt-work-and-codex}
    {OpenAI Work and Codex documentation}
\end{itemize}

% ============================================================
\section*{Final Operating Principle}
\addcontentsline{toc}{section}{Final Operating Principle}
% ============================================================

\begin{tcolorbox}[
    enhanced,
    breakable,
    colback=CUNavy,
    colframe=CUGold,
    coltitle=CUNavy,
    colbacktitle=CUGoldLight,
    boxrule=1pt,
    arc=2mm,
    title=\textbf{Controlled Development Rule}
]
\color{white}

Plan one feature.

Implement one feature.

Review every changed file.

Test locally.

Run the production build.

Commit only to \branchname{website}.

Push only after the change is understood.

Do not deploy until the local website is stable.

\end{tcolorbox}

\vfill

\begin{center}
\color{CUMuted}
\small
End of \DocumentTitle
\end{center}

\end{document}
